\begin{name}
    {\tenchude}
    {\tendethi}
    {Toán thầy Phát}
    {\thoigian}
\end{name}
\Opensolutionfile{ans}[ans/ans-2-GHK1-8-GHK1-LuongTheVinh-CSA-HaNoi-20]
\setcounter{ex}{0}\setcounter{bt}{0}
\begin{ex}%[GHK1 - Lương Thế Vinh CSA]%[Phan Anh - EX1]%[2D1Y1-1]
	Cho hàm số $y=\dfrac{1}{3}x^3-2x^2+3x+1$. Khẳng định nào sau đây đúng?
	\choice
	{\True Hàm số nghịch biến trên khoảng $\left(1;3\right)$}
	{Hàm số nghịch biến trên khoảng $\left(-\infty;1\right)$}
	{Hàm số đồng biến trên khoảng $\left(1;3\right)$}
	{Hàm số nghịch biến trên khoảng $(3;+\infty)$}
	\loigiai{\immini{Ta có $y'=x^2-4x+3$.\\
			$y'=0\Leftrightarrow x^2-4x+3=0\Leftrightarrow\hoac{&x=1\\&x=3.}$\\
			Bảng biến thiên của hàm số như hình vẽ bên.\\
			Suy ra hàm số nghịch biến trên $(1;3)$.}
		{\begin{tikzpicture}
			\tkzTabInit[nocadre=false,lgt=1.2,espcl=1.6,deltacl=0.6]
			{$x$/0.8,$y'$/0.8,$y$/2.2}{$-\infty$,$1$,$3$,$+\infty$}
			\tkzTabLine{,+,0,-,0,+,}
			\tkzTabVar{-/$ $,+/$ $,-/$ $,+/$ $}	
			\end{tikzpicture}}
	}
\end{ex}

\begin{ex}%[GHK1 - Lương Thế Vinh CSA]%[Phan Anh - EX1]%[2D1B1-1]
	Hàm số nào dưới đây đồng biến trên $(-\infty;+\infty)$?
	\choice
	{\True $y=x^3-x^2+2019x$}
	{$y=\tan x$}
	{$y=\dfrac{x+2}{x+2019}$}
	{$y=x^4+2019$}
	\loigiai{Xét từng đáp án ta được
		\begin{itemize}
			\item Hàm số $y=x^3-x^2+2019x$ có $y'=3x^2-2x+2019>0$, $\forall x\in\mathbb{R}$ nên đồng biến trên $\mathbb{R}$.
			\item Hàm số $y=\tan x$ có tập xác định là $\mathscr{D}=\mathbb{R}\setminus\left\{\dfrac{\pi}{2}+k\pi,k\in\mathbb{Z}\right\}$, nên không đồng biến trên $\mathbb{R}$.
			\item Hàm số $y=\dfrac{x+2}{x+2019}$ có tập xác định là $\mathscr{D}=\mathbb{R}\setminus\{-2019\}$, nên không đồng biến trên $\mathbb{R}$.
			\item Hàm số $y=x^4+2019$ có $y'=4x^3$ đổi dấu khi đi qua giá trị $x=0$, nên không đồng biến trên $\mathbb{R}$.
	\end{itemize}}
\end{ex}

\begin{ex}%[GHK1 - Lương Thế Vinh CSA]%[Phan Anh - EX1]%[2D1B2-1]
	Hàm số nào sau đây có cực trị?
	\choice
	{\True $y=x^4-2x^2+2019$}
	{$y=x^3-3x^2+7x$}
	{$y=\tan x$}
	{$y=\dfrac{x+1}{x+2}$}
	\loigiai{\immini{Hàm số $y=x^4-2x^2+2019$ có $y'=4x^3-4x$.\\
			$y'=0\Leftrightarrow\hoac{&x=0\\&x=1\\&x=-1.}$\\
			Ta có bảng biến thiên như hình vẽ bên.
			Vậy hàm số $y=x^4-2x^2+2019$ có ba điểm cực trị.}
		{\begin{tikzpicture}
			\tkzTabInit[nocadre=false,lgt=1.2,espcl=1.6,deltacl=0.6]
			{$x$/0.8,$f'(x)$/0.8,$f(x)$/2.5}{$-\infty$,$-1$,$0$,$1$,$+\infty$}
			\tkzTabLine{,-,0,+,0,-,0,+,}
			\tkzTabVar{+/,-/,+/,-/,+/}	
			\end{tikzpicture}}}
\end{ex}

\begin{ex}%[GHK1 - Lương Thế Vinh CSA]%[Phan Anh - EX1]%[2D1B2-1] Hàm số $y=x^3-3x^2+2019$ đạt cực đại tại điểm nào?
	\choice
	{\True $x=0$}
	{$x=1$}
	{$x=2$}
	{$x=2019$}
	\loigiai{\immini{Ta có $y'=3x^2-6x$, $y'=0\Leftrightarrow\hoac{&x=0\\&x=2.}$\\
			Ta có bảng biến thiên như hình vẽ bên.\\
			Từ bảng biến thiên ta suy ra hàm số đạt cực đại tại điểm $x=0$.}
		{\begin{tikzpicture}
			\tkzTabInit[nocadre=false,lgt=1.2,espcl=1.6,deltacl=0.6]
			{$x$/0.8,$f'(x)$/0.8,$f(x)$/2.5}{$-\infty$,$0$,$2$,$+\infty$}
			\tkzTabLine{,+,0,-,0,+,}
			\tkzTabVar{-/,+/CĐ,-/CT,+/}	
			\end{tikzpicture}}}
\end{ex}

\begin{ex}%[GHK1 - Lương Thế Vinh CSA]%[Phan Anh - EX1]%[2D1B2-1]
	Đồ thị hàm số $y=x^3-3x^2-1$ có hai điểm cực trị là $A$, $B$. Tính độ dài đoạn thẳng $AB$. 
	\choice
	{\True $2\sqrt{5}$}
	{$2$}
	{$3\sqrt{2}$}
	{$2\sqrt{3}$}
	\loigiai{Ta có $y'=3x^2-6x$, $y'=0\Leftrightarrow\hoac{&x=0&\Rightarrow&y=-1\\&x=2&\Rightarrow&y=-5.}$\\
		Vậy đồ thị hàm số có hai điểm cực trị là $A(0;-1)$ và $B(2;-5)$.\\
		Suy ra $AB=\sqrt{(2-0)^2+(-5+1)^2}=2\sqrt{5}$.}
\end{ex}

\begin{ex}%[GHK1 - Lương Thế Vinh CSA]%[Phan Anh - EX1]%[2D1Y4-1]
	Cho hàm số $y=\dfrac{x+1}{2x-2}$. Khẳng định nào sau đây đúng?
	\choice
	{\True Đồ thị hàm số có tiệm cận ngang là $y=\dfrac{1}{2}$}
	{Đồ thị hàm số có tiệm cận ngang là $x=\dfrac{1}{2}$}
	{Đồ thị hàm số có tiệm cận đứng là $x=\dfrac{1}{2}$}
	{Đồ thị hàm số có tiệm cận đứng là $x=2$}
	\loigiai{Ta có $\lim\limits_{x\to+\infty}y=\lim\limits_{x\to+\infty}\dfrac{x+1}{2x-2}=\dfrac{1}{2}$ và $\lim\limits_{x\to-\infty}y=\lim\limits_{x\to-\infty}\dfrac{x+1}{2x-2}=\dfrac{1}{2}$ nên đồ thị hàm số nhận $y=\dfrac{1}{2}$ làm tiệm cận ngang.}
\end{ex}

\begin{ex}%[GHK1 - Lương Thế Vinh CSA]%[Phan Anh - EX1]%[2D1B4-1]
	Đồ thị hàm số nào dưới đây có tiệm cận ngang?
	\choice
	{\True $y=\dfrac{\sqrt{x-2019}}{x+2019}$}
	{$y=x^3-x$}
	{$y=\dfrac{x+\sqrt{4-x^2}}{x}$}
	{$y=\dfrac{x^2+2019}{x}$}
	\loigiai{Xét tứng đáp án ta có
		\begin{itemize}
			\item Hàm số $y=\dfrac{\sqrt{x-2019}}{x+2019}$ có $\lim\limits_{x\to+\infty}y=0$ nên đồ thị nhận $y=0$ làm tiệm cận ngang.
			\item Hàm số $y=x^3-x$ có $\lim\limits_{x\to+\infty}y=+\infty$ và $\lim\limits_{x\to-\infty}y=-\infty$ nên đồ thị hàm số không có tiệm cận ngang.
			\item Hàm số $y=\dfrac{x+\sqrt{4-x^2}}{x}$ có tập xác định là $\mathscr{D}=[-2;2]\setminus\{0\}$ nên đồ thị hàm số không có tiệm cận ngang.
			\item Hàm số $y=\dfrac{x^2+2019}{x}$ có $\lim\limits_{x\to+\infty}y=+\infty$ và $\lim\limits_{x\to-\infty}y=-\infty$ nên đồ thị hàm số không có tiệm cận ngang.
	\end{itemize}}
\end{ex}

\begin{ex}%[GHK1 - Lương Thế Vinh CSA]%[Phan Anh - EX1]%[2D1B5-3]
	\immini{Cho hàm số $y=f(x)$ có bảng biến thiên như hình bên. Phương trình $2f(x)-5=0$ có bao nhiêu nghiệm?
		\choice
		{\True $4$}
		{$1$}
		{$2$}
		{$3$}}
	{\begin{tikzpicture}
		\tkzTabInit[nocadre=false,lgt=1.2,espcl=1.8,deltacl=0.6]{$x$/0.7,$y'$/0.7,$y$/2}{$-\infty$,$-1$,$0$,$1$,$+\infty$}
		\tkzTabLine{,+,$0$,-,$0$,+,$0$,-}
		\tkzTabVar{-/$-\infty$,+/$3$,-/$1$,+/$3$,-/$-\infty$}
		\end{tikzpicture}}
	\loigiai{Ta có $2f(x)-5=0\Leftrightarrow f(x)=\dfrac{5}{2}$.\\
		Từ bảng biến thiên ta suy ra phương trình có $4$ nghiệm phân biệt.}
\end{ex}

\begin{ex}%[GHK1 - Lương Thế Vinh CSA]%[Phan Anh - EX1]%[2D1B5-4]
	Đồ thị hàm số $y=-x^4+3x^2$ và trục hoành có bao nhiêu giao điểm?
	\choice
	{\True $3$}
	{$1$}
	{$2$}
	{$0$}
	\loigiai{Phương trình hoành độ giao điểm của đồ thị với trục hoành là
		$$-x^4+3x^2=0\Leftrightarrow\hoac{&x^2=0\\&x^2=3}\Leftrightarrow\hoac{&x=0\\&x=\sqrt{3}\\&x=-\sqrt{3}.}$$
		Vậy đồ thị hàm số cắt trục hoành tại ba điểm phân biệt.}
\end{ex}

\begin{ex}%[GHK1 - Lương Thế Vinh CSA]%[Phan Anh - EX1]%[2D1Y3-1]
	Giá trị lớn nhất của hàm số $y=x^4-4x^2-5$ trên đoạn $[-2;3]$ bằng
	\choice
	{\True $40$}
	{$-5$}
	{$-9$}
	{$112$}
	\loigiai{Ta có $y'=4x^3-8x$, $y'=0\Leftrightarrow\hoac{&x=0\\&x=\sqrt{2}\\&x=-\sqrt{2}.}$\\
		Xét trên $[-2;3]$, ta có $f(0)=-5$, $f(-2)=-5$, $f(3)=40$, $f\left(-\sqrt{2}\right)=f\left(\sqrt{2}\right)=-9$.\\
		Vậy giá trị lớn nhất của hàm số trên $[-2;3]$ là $40$.}
\end{ex}

\begin{ex}%[GHK1 - Lương Thế Vinh CSA]%[Phan Anh - EX1]%[2D1Y5-1]
	\immini{Biết rằng đường cong trong hình bên là đồ thị của một trong bốn hàm số được liệt kê ở bốn phương án dưới đây. Hỏi hàm số đó là hàm số nào?
		\choice
		{\True $y=x^3-3x+1$}
		{$y=-x^3-3x+1$}
		{$y=x^4-x^2+3$}
		{$y=x^2-3x+1$}}
	{\begin{tikzpicture}[scale=0.7, font=\footnotesize, line join=round, line cap=round, >=stealth]
		\draw[->](-3,0)--(3,0)node[above]{$x$};
		\draw[->](0,-2)--(0,4)node[right]{$y$};
		\draw[smooth,samples=100,domain=-2.1:2.1]plot(\x,{(\x)^3-3*(\x)+1});
		\fill (0,0)node[below left]{$O$}circle(1.2pt);
		\end{tikzpicture}}
	\loigiai{Dựa vào hình dáng đồ thị, ta nhận thấy đó là đồ thị hàm số bậc ba với hệ số $a>0$ nên nhận hàm số $y=x^3-3x+1$.}
\end{ex}

\begin{ex}%[GHK1 - Lương Thế Vinh CSA]%[Phan Anh - EX1]%[1D5B2-2]
	Viết phương trình tiếp tuyến của đồ thị hàm số $y=x^3-3x^2-1$ tại điểm $M(3;-1)$. 
	\choice
	{\True $y=9x-28$}
	{$y=3x-10$}
	{$y=x-4$}
	{$y=-4x+2$}
	\loigiai{Ta có $y'=3x^2-6x$, $y'(3)=9$. Phương trình tiếp tuyến của đồ thị tại $M$ là
		$$y=y'(3)(x-3)-1\Leftrightarrow y=9(x-3)-1\Leftrightarrow y=9x-28.$$}
\end{ex}

\begin{ex}%[GHK1 - Lương Thế Vinh CSA]%[Phan Anh - EX1]%[2D2B1-2]
	Rút gọn biểu thức $P=\dfrac{\sqrt[6]{x}\sqrt[3]{x^4}\sqrt[4]{x}}{\sqrt{x^3}}$ với $x$ là số thực dương.
	\choice
	{\True $x^{\tfrac{1}{4}}$}
	{$x^{\tfrac{13}{15}}$}
	{$x^{\tfrac{13}{18}}$}
	{$x^{\tfrac{1}{6}}$}
	\loigiai{Ta có $P=\dfrac{\sqrt[6]{x}\sqrt[3]{x^4}\sqrt[4]{x}}{\sqrt{x^3}}=x^{\tfrac{1}{6}+\tfrac{4}{3}+\tfrac{1}{4}-\tfrac{3}{2}}=x^{\tfrac{1}{4}}$.}
\end{ex}

\begin{ex}%[GHK1 - Lương Thế Vinh CSA]%[Phan Anh - EX1]%[2D2B3-2]
	Với $a$ là số thực khác $0$, khẳng định nào dưới đây đúng?
	\choice
	{\True $\log(a^4)=4\log|a|$}
	{$\log(a^4)=\dfrac{1}{4}\log a$}
	{$\log(a^4)=4\log a$}
	{$\log(4a)=\dfrac{1}{4}\log a$}
	\loigiai{Ta có khẳng định đúng là ``$\log(a^4)=4\log|a|$''.}
\end{ex}

\begin{ex}%[GHK1 - Lương Thế Vinh CSA]%[Phan Anh - EX1]%[2D2Y3-2]
	Cho $a$, $b$ là các số thực dương tùy ý, khẳng định nào dưới đây đúng?
	\choice
	{\True $\log(ab)=\log a+\log b$}
	{$\log(ab)=\log a\log b$}
	{$\log(a+b)=\log a\log b$}
	{$\log(a+b)=\log a+\log b$}
	\loigiai{Ta có công thức đúng là $\log(ab)=\log a+\log b$.}
\end{ex}

\begin{ex}%[GHK1 - Lương Thế Vinh CSA]%[Phan Anh - EX1]%[2D1K4-1]
	Đồ thị hàm số $y=\dfrac{\sqrt{x+3}}{(x-2)^3(x+3)}$ có bao nhiêu tiệm cận đứng?
	\choice
	{\True $2$}
	{$1$}
	{$4$}
	{$0$}
	\loigiai{Tập xác định của hàm số là $\mathscr{D}=(-3;+\infty)\setminus\{2\}$.\\
		Ta có $y=\dfrac{\sqrt{x+3}}{(x-2)^3(x+3)}=\dfrac{1}{(x-2)^3\sqrt{x+3}}$.
		\begin{itemize}
			\item $\lim\limits_{x\to(-3)^+}y=-\infty$ nên đồ thị hàm số nhận $x=-3$ làm tiệm cận đứng.
			\item $\lim\limits_{x\to(2)^+}y=+\infty$ và $\lim\limits_{x\to(2)^-}y=-\infty$ nên đồ thị hàm số nhận $x=2$ làm tiệm cận đứng.
		\end{itemize}
		Vậy đồ thị hàm số có $2$ đường tiệm cận đứng.}
\end{ex}

\begin{ex}%[GHK1 - Lương Thế Vinh CSA]%[Phan Anh - EX1]%[2H1B3-2]
	Cho hình chóp $SABC$ có $SA$, $SB$, $SC$ đôi một vuông góc và $SA=a$, $SB=3a$, $SC=4a$. Thể tích khối chóp $SABC$ tính theo $a$ là
	\choice
	{\True $2a^3$}
	{$a^3$}
	{$4a^3$}
	{$12a^3$}
	\loigiai{Do $S.ABC$ có $SA$, $SB$, $SC$ đôi một vuông góc nên $V=\dfrac{1}{6}\cdot SA\cdot SB\cdot SC=2a^3$.}
\end{ex}

\begin{ex}%[GHK1 - Lương Thế Vinh CSA]%[Phan Anh - EX1]%[2H1B3-2]
	Cho hình chóp $SABC$ có đáy $ABC$ là tam giác vuông tại $B$, $AB=a$ và $AC=a\sqrt{3}$. Biết $SA\perp (ABC)$ và $SB=a\sqrt{5}$. Tính theo $a$ thể tích khối chóp $SABC$.
	\choice
	{\True $\dfrac{a^3\sqrt{2}}{3}$}
	{$\dfrac{a^3\sqrt{6}}{6}$}
	{$\dfrac{a^3\sqrt{15}}{6}$}
	{$\dfrac{a^3\sqrt{6}}{4}$}
	\loigiai{\immini{Xét tam giác $ABC$ ta có $BC=\sqrt{AC^2-AB^2}=a\sqrt{2}$.\\
			Xét tam giác $SAB$ ta có $SA=\sqrt{SB^2-AB^2}=2a$.\\
			Thể tích của khối chóp $S.ABC$ là
			$$V=\dfrac{1}{3}\cdot SA\cdot S_{ABC}=\dfrac{1}{3}\cdot2a\cdot\dfrac{1}{2}\cdot a\cdot a\sqrt{2}=\dfrac{a^3\sqrt{2}}{3}.$$}
		{\begin{tikzpicture}[line cap=round, line join=round,font=\footnotesize,>=stealth, scale=0.5]
			\tkzDefPoints{0/0/A,2/-1.5/B,6/0/C}
			\coordinate (S) at ($(A)+(0,6)$);
			\tkzDrawPoints[fill=black](A,B,C,S)
			\tkzDrawSegments(S,A S,B S,C A,B B,C)
			\tkzDrawSegments[dashed](A,C)
			\tkzLabelPoints[above](S)
			\tkzLabelPoints[above left](A)
			\tkzLabelPoints[above right](C)
			\tkzLabelPoints[below](B)
			\tkzMarkRightAngle(C,B,A)
			\tkzMarkRightAngle(B,A,S)
			\end{tikzpicture}}}
\end{ex}

\begin{ex}%[GHK1 - Lương Thế Vinh CSA]%[Phan Anh - EX1]%[2H1B3-2]
	Cho hình chóp tứ giác $S.ABCD$ có đáy là hình chữ nhật, $AB=2a,BC=a$. Mặt bên $(SAB)$ là tam giác đều nằm trong mặt phẳng vuông góc với mặt đáy. Thể tích khối chóp $S.ABCD$ là
	\choice
	{\True $\dfrac{2a^3\sqrt{3}}{3}$}
	{$\dfrac{a^3\sqrt{3}}{6}$}
	{$\dfrac{a^3\sqrt{3}}{2}$}
	{$\dfrac{a^3\sqrt{3}}{3}$}
	\loigiai{\immini{Gọi $H$ là trung điểm của $AB$, ta có $$\heva{&SH\perp AB\\&(SAB)\perp(ABCD)\\&(SAB)\cap(ABCD)=AB}\Rightarrow SH\perp(ABCD).$$
			Tam giác $SAB$ đều nên $SH=\dfrac{AB\sqrt{3}}{2}=a\sqrt{3}$.\\
			Thể tích của khối chóp là 
			$$V=\dfrac{1}{3}\cdot SH\cdot S_{ABCD}=\dfrac{1}{3}\cdot a\sqrt{3}\cdot2a\cdot a=\dfrac{2a^3\sqrt{3}}{3}.$$}
		{\begin{tikzpicture}[line cap=round, line join=round,font=\footnotesize,>=stealth, scale=0.7]
			\tkzDefPoints{0/0/A,5/0/D,-2/-1.2/B}
			\coordinate (C) at ($(B)+(D)-(A)$);
			\coordinate (H) at ($(B)!0.5!(A)$);
			\coordinate (S) at ($(H)+(0,6)$);
			\tkzDrawPoints[fill=black](A,B,C,D,S,H)
			\tkzDrawSegments(S,B S,C S,D B,C C,D)
			\tkzDrawSegments[dashed](A,B S,A A,D S,H)
			\tkzLabelPoints[above](S)
			\tkzLabelPoints[above left](A,H)
			\tkzLabelPoints[below left](B)
			\tkzLabelPoints[below right](C)
			\tkzLabelPoints[above right](D)
			\end{tikzpicture}}}
\end{ex}

\begin{ex}%[GHK1 - Lương Thế Vinh CSA]%[Phan Anh - EX1]%[2H1B3-2]
	Cho hình chóp tứ giác đều có tất cả các cạnh bằng $a$. Tính thể tích của khối chóp đã cho
	\choice
	{\True $\dfrac{\sqrt{2}a^3}{6}$}
	{$\dfrac{4a^3}{3}$}
	{$\dfrac{8\sqrt{2}a^3}{3}$}
	{$V=\dfrac{2\sqrt{2}a^3}{3}$}
	\loigiai{\immini{Gọi $O$ là tâm của hình vuông $ABCD$. Ta có $SO\perp(ABCD)$.\\
			Có $AO=\dfrac{1}{2}AC=\dfrac{a\sqrt{2}}{2}$.\\
			Xét tam giác $SAO$ ta có $SO=\sqrt{SA^2-AO^2}=\dfrac{a\sqrt{2}}{2}$.\\
			Thể tích khối chóp là $$V=\dfrac{1}{3}\cdot SO\cdot S_{ABCD}=\dfrac{1}{3}\cdot\dfrac{a\sqrt{2}}{2}\cdot a^2=\dfrac{a^3\sqrt{2}}{6}.$$}
		{\begin{tikzpicture}[line cap=round, line join=round,font=\footnotesize,>=stealth, scale=0.9]
			\tkzDefPoints{0/0/A,5/0/D,-2/-1.2/B}
			\coordinate (C) at ($(B)+(D)-(A)$);
			\coordinate (O) at ($(B)!0.5!(D)$);
			\coordinate (S) at ($(O)+(0,6)$);
			\tkzDrawPoints[fill=black](A,B,C,D,S,O)
			\tkzDrawSegments(S,B S,C S,D B,C C,D)
			\tkzDrawSegments[dashed](A,B S,A A,D S,O A,C B,D)
			\tkzLabelPoints[above](S)
			\tkzLabelPoints[above left](A)
			\tkzLabelPoints[below left](B)
			\tkzLabelPoints[below right](C)
			\tkzLabelPoints[above right](D,O)
			\end{tikzpicture}}}
\end{ex}

\begin{ex}%[GHK1 - Lương Thế Vinh CSA]%[Phan Anh - EX1]%[2H2B2-2]
	Cho hình chóp $S.ABC$ có đáy $ABC$ là tam giác vuông tại $A$, $AB=2$, $AC=4$. $SA$ vuông góc với mặt phẳng $(ABC)$ và $SA=\sqrt{5}$. Mặt cầu đi qua các đỉnh hình chóp $S.ABC$ có bán kính $R$ bằng bao nhiêu?
	\choice
	{\True $R=\dfrac{5}{2}$}
	{$R=\dfrac{10}{3}$}
	{$R=5$}
	{$R=\dfrac{25}{2}$}
	\loigiai{\immini{Gọi $O$ là trung điểm của $BC$, kẻ $Ox\parallel SA$, ta có $Ox$ là trục đường tròn ngoại tiếp tam giác $ABC$.\\
			Gọi $M$ là trung điểm của $SA$, kẻ trung trực $Mt$ của $SA$ trong mặt phẳng $(SAO)$, cắt $Ox$ tại $I$. Ta có $I$ là tâm mặt cầu ngoại tiếp $S.ABC$.\\
			Ta có bán kính của mặt cầu là
			$$R=IA=\sqrt{AM^2+AO^2}=\dfrac{1}{2}\sqrt{SA^2+AB^2+AC^2}=\dfrac{5}{2}.$$}
		{\begin{tikzpicture}[line cap=round, line join=round,font=\footnotesize,>=stealth, scale=0.6]
			\tkzDefPoints{0/0/A,-2/-1.5/B,6/0/C}
			\coordinate (S) at ($(A)+(0,6)$);
			\coordinate (M) at ($(S)!0.5!(A)$);
			\coordinate (O) at ($(B)!0.5!(C)$);
			\coordinate (I) at ($(M)+(O)-(A)$);
			\coordinate (t) at ($(M)!1.5!(I)$);
			\coordinate (x) at ($(O)!1.9!(I)$);
			\tkzInterLL(M,I)(S,O)	\tkzGetPoint{i}
			\tkzInterLL(A,I)(S,O)	\tkzGetPoint{a}
			\tkzDrawPoints[fill=black](A,B,C,S,M,O,I)
			\tkzDrawSegments(S,B S,C B,C O,x i,t S,O I,a)
			\tkzDrawSegments[dashed](S,A A,B A,C M,i A,O A,a)
			\tkzLabelPoints[above](S)
			\tkzLabelPoints[above left](A,M)
			\tkzLabelPoints[above right](C,x,I)
			\tkzLabelPoints[below](B,O,t)
			\end{tikzpicture}}}
\end{ex}
\begin{ex}%[GHK1 - Lương Thế Vinh CSA]%[Phan Anh - EX1]%[2D1B2-3]
	Tìm tất cả các giá trị của tham số $m$ để hàm số $f(x)=\dfrac{1}{3}x^3+mx^2+(m^2-4)x$ đạt cực tiểu tại $x=1$.
	\choice
	{\True $m=1$}
	{$m=3$}
	{$m\in\{1;-3\}$}
	{$m=-3$}
	\loigiai{Ta có $f'(x)=x^2+2mx+m^2-4$, $f''(x)=2x+2m$. Hàm số đạt cực tiểu tại $x=1$ khi và chỉ khi
		$$\heva{&f'(1)=0\\&f''(1)>0}\Leftrightarrow\heva{&m^2+2m-3=0\\&2+2m>0}\Leftrightarrow\heva{&m=1\vee m=-3\\&m>-1}\Leftrightarrow m=1.$$}
\end{ex}

\begin{ex}%[GHK1 - Lương Thế Vinh CSA]%[Phan Anh - EX1]%[2D1K3-1]
	Giá trị lớn nhất của hàm số $y=2\cos^2 x+\sin x+1$ là phân số tối giản có dạng $\dfrac{a}{b}$ với $a,b$ là các số nguyên dương. Tìm $a+b$.
	\choice
	{\True $33$}
	{$17$}
	{$9$}
	{$10$}
	\loigiai{Ta có $y=2\cos^2 x+\sin x+1=-2\sin^2x+\sin x+3$, đặt $t=\sin x$, ta được
		$$y=f(t)=-2t^2+t+3,\quad t\in[-1;1].$$
		Ta có $f'(t)=-4t+1$, $f'(t)=0\Leftrightarrow t=\dfrac{1}{4}$.\\
		Xét trên đoạn $[-1;1]$, ta có $f(-1)=1$, $f(1)=2$, $f\left(\dfrac{1}{4}\right)=\dfrac{25}{8}$.\\
		Vậy $\max y=\max\limits_{t\in[-1;1]}f(t)=\dfrac{25}{8}$, suy ra $a=25$ và $b=8$ nên $a+b=33$.}
\end{ex}

\begin{ex}%[GHK1 - Lương Thế Vinh CSA]%[Phan Anh - EX1]%[2D2B3-2]
	Cho $\log_43=a$ và $\log_73=b$. Mệnh đề nào sau đây đúng?
	\choice
	{\True $\log_{28}3=\dfrac{ab}{a+b}$}
	{$\log_{28}3=\dfrac{1}{a+b}$}
	{$\log_{28}3=\dfrac{a+b}{ab}$}
	{$\log_{28}3=\dfrac{1}{ab}$}
	\loigiai{Ta có $\log_{28}3=\dfrac{1}{\log_328}=\dfrac{1}{\log_34+\log_37}=\dfrac{1}{\dfrac{1}{a}+\dfrac{1}{b}}=\dfrac{ab}{a+b}$.}
\end{ex}

\begin{ex}%[GHK1 - Lương Thế Vinh CSA]%[Phan Anh - EX1]%[2D2B3-2]
	Với mọi số thực dương $x$, $y$ thỏa mãn $x^2+y^2=8xy$, mệnh đề nào dưới đây đúng? 
	\choice
	{\True $\log(x+y)=\dfrac{1}{2}(1+\log x+\log y)$}
	{$\log(x+y)=\dfrac{1}{2}(\log x+\log y)$}
	{$\log(x+y)=1+\log x+\log y$}
	{$\log(x+y)=\dfrac{1}{2}+\log x+\log y$}
	\loigiai{Ta có $x^2+y^2=8xy\Leftrightarrow(x+y)^2=10xy$. Lấy logarit thập phân hai vế ta được
		\begin{eqnarray*}
			\log(x+y)^2=\log(10xy)&\Leftrightarrow&2\log(x+y)=1+\log x+\log y\\
			&\Leftrightarrow&\log(x+y)=\dfrac{1}{2}(1+\log x+\log y).
	\end{eqnarray*}}
\end{ex}

\begin{ex}%[GHK1 - Lương Thế Vinh CSA]%[Phan Anh - EX1]%[2D2B3-2]
	Cho các số thực dương $x$, $y$ thỏa mãn $\log(xy^2)=5$ và $\log(x^3y)=10$. Tính $P=\log(xy)$.
	\choice
	{\True $P=4$}
	{$P=2$}
	{$P=5$}
	{$P=1$}
	\loigiai{Từ giả thiết ta có $\heva{&\log x+2\log y=5\\&3\log x+\log y=10}\Leftrightarrow\heva{&\log x=3\\&\log y=1.}$\\
		Vậy ta có $P=\log xy=\log x+\log y=4$.}
\end{ex}

\begin{ex}%[GHK1 - Lương Thế Vinh CSA]%[Phan Anh - EX1]%[2D2K1-2]
	Tính giá trị biểu thức $P=\left(17-12\sqrt{2}\right)^{2017}\cdot\left(7+5\sqrt{2}\right)^{2690}$.
	\choice
	{\True $P=3+2\sqrt{2}$}
	{$P=\sqrt{2}+1$}
	{$P=\sqrt{2}-1$}
	{$P=3-2\sqrt{2}$}
	\loigiai{Ta có 
		\begin{eqnarray*}
			P&=&\left(17-12\sqrt{2}\right)^{2017}\cdot\left(7+5\sqrt{2}\right)^{2690}\\
			&=&\left(\left(\sqrt{2}-1\right)^4\right)^{2017}\cdot\left(\left(\sqrt{2}+1\right)^3\right)^{2690}\\
			&=&\left(\sqrt{2}-1\right)^{8068}\cdot\left(\sqrt{2}+1\right)^{8070}\\
			&=&\left(\left(\sqrt{2}-1\right)\cdot\left(\sqrt{2}+1\right)\right)^{8068}\cdot\left(\sqrt{2}+1\right)^2=3+2\sqrt{2}.
	\end{eqnarray*}}
\end{ex}

\begin{ex}%[GHK1 - Lương Thế Vinh CSA]%[Phan Anh - EX1]%[2D2K3-2]
	Cho $\log_23=x$ và $\log_25=y$. Biết rằng $\log_{20}15=\dfrac{ax+by}{cy+2}$ với $a$, $b$, $c$ là các số nguyên dương. Tính $P=a+b+c$.
	\choice
	{\True $P=3$}
	{$P=4$}
	{$P=5$}
	{$P=6$}
	\loigiai{Ta có $\log_{20}15=\dfrac{\log_215}{\log_220}=\dfrac{\log_23+\log_25}{\log_24+\log_25}=\dfrac{x+y}{2+y}$.\\
		Vậy nên $a=b=c=1$, suy ra $P=a+b+c=3$.}
\end{ex}

\begin{ex}%[GHK1 - Lương Thế Vinh CSA]%[Phan Anh - EX1]%[2H1B3-2]
	Cho hình chóp $S.ABCD$ có đáy $ABCD$ là hình vuông , $SA\perp (ABCD)$, $SA=2a\sqrt{3}$, $SD$ tạo với đáy một góc $60^\circ$. Thể tích của khối chóp $S.ABCD$ là
	\choice
	{\True $\dfrac{8a^3\sqrt{3}}{3}$}
	{$\dfrac{4a^3\sqrt{3}}{3}$}
	{$\dfrac{2a^3\sqrt{3}}{3}$}
	{$a^3\sqrt{3}$}
	\loigiai{\immini{Ta có $SA\perp(ABCD)$ nên $AD$ là hình chiếu của $SD$ lên mặt đáy $(ABCD)$, suy ra góc giữa $SD$ với đáy là $\widehat{SDA}=60^\circ$.\\
			Tam giác $SAD$ có $AD=\dfrac{SA}{\tan60^\circ}=2a$.\\
			Thể tích của khối chóp $S.ABCD$ là
			$$V=\dfrac{1}{3}\cdot SA\cdot S_{ABCD}=\dfrac{1}{3}\cdot2a\sqrt{3}\cdot4a^2=\dfrac{8a^3\sqrt{3}}{3}.$$}
		{\begin{tikzpicture}[line cap=round, line join=round,font=\footnotesize,>=stealth, scale=0.6]
			\tkzDefPoints{0/0/A,5/0/D,-2/-1.2/B}
			\coordinate (C) at ($(B)+(D)-(A)$);
			\coordinate (S) at ($(A)+(0,6)$);
			\tkzDrawPoints[fill=black](A,B,C,D,S)
			\tkzDrawSegments(S,B S,C S,D B,C C,D)
			\tkzDrawSegments[dashed](A,B S,A A,D)
			\tkzLabelPoints[above](S)
			\tkzLabelPoints[above left](A)
			\tkzLabelPoints[below left](B)
			\tkzLabelPoints[below right](C)
			\tkzLabelPoints[above right](D)
			\end{tikzpicture}}}
\end{ex}

\begin{ex}%[GHK1 - Lương Thế Vinh CSA]%[Phan Anh - EX1]%[2H1B3-2]
	Cho hình chóp $S.ABCD$ có đáy $ABCD$ là hình vuông cạnh $a$, $SD=\dfrac{3a}{2}$, hình chiếu vuông góc của $S$ trên mặt phẳng $(ABCD)$ là trung điểm cạnh $AB$. Tính theo $a$ thể tích khối chóp $S.ABCD$
	\choice
	{\True $\dfrac{a^3}{3}$}
	{$\dfrac{a^3\sqrt{3}}{3}$}
	{$\dfrac{a^3}{6}$}
	{$\dfrac{a^3\sqrt{3}}{2}$}
	\loigiai{\immini{Xét tam giác $AHD$, ta có $HD=\sqrt{AH^2+AD^2}=\dfrac{a\sqrt{5}}{2}$.\\
			Xét tam giác $SHD$, ta có $SH=\sqrt{SD^2-HD^2}=a$.\\
			Thể tích của khối chóp là $$V=\dfrac{1}{3}\cdot SH\cdot S_{ABCD}=\dfrac{1}{3}\cdot a\cdot a^2=\dfrac{a^3}{3}.$$}
		{\begin{tikzpicture}[line cap=round, line join=round,font=\footnotesize,>=stealth, scale=0.7]
			\tkzDefPoints{0/0/A,5/0/D,-2/-1.2/B}
			\coordinate (C) at ($(B)+(D)-(A)$);
			\coordinate (H) at ($(B)!0.5!(A)$);
			\coordinate (S) at ($(H)+(0,6)$);
			\tkzDrawPoints[fill=black](A,B,C,D,S,H)
			\tkzDrawSegments(S,B S,C S,D B,C C,D)
			\tkzDrawSegments[dashed](A,B S,A A,D S,H H,D)
			\tkzLabelPoints[above](S)
			\tkzLabelPoints[above left](A,H)
			\tkzLabelPoints[below left](B)
			\tkzLabelPoints[below right](C)
			\tkzLabelPoints[above right](D)
			\end{tikzpicture}}}
\end{ex}

\begin{ex}%[GHK1 - Lương Thế Vinh CSA]%[Phan Anh - EX1]%[2H1B3-2]
	Cho hình chóp tam giác đều $S.ABC$ có độ dài cạnh đáy bằng $a$, góc giữa cạnh bên và mặt đáy bằng $60^\circ$. Tính thể tích khối chóp đã cho
	\choice
	{\True $\dfrac{\sqrt{3}a^3}{12}$}
	{$\dfrac{\sqrt{3}a^3}{6}$}
	{$\dfrac{\sqrt{3}a^3}{3}$}
	{$\dfrac{\sqrt{3}a^3}{4}$}
	\loigiai{\immini{Gọi $O$ là tâm tam giác đều $ABC$, ta có $AO$ là hình chiếu của $SA$ lên mặt đáy $ABC$, suy ra $\widehat{SAO}=60^\circ$.\\
			Gọi $M$ là trung điểm của $BC$, ta có $$AO=\dfrac{2}{3}AM=\dfrac{2}{3}\cdot\dfrac{a\sqrt{3}}{2}=\dfrac{a\sqrt{3}}{3}.$$
			Tam giác $SAO$ có $SO=AO\tan60^\circ=a$.\\
			Thể tích khối chóp là
			$$V=\dfrac{1}{3}\cdot SO\cdot S_{ABC}=\dfrac{1}{3}\cdot a\cdot\dfrac{a^2\sqrt{3}}{4}=\dfrac{a^3\sqrt{3}}{12}.$$}
		{\begin{tikzpicture}[line cap=round, line join=round,font=\footnotesize,>=stealth, scale=0.8]
			\tkzDefPoints{0/0/A,2/-1.5/B,6/0/C}
			\coordinate (M) at ($(B)!0.5!(C)$);
			\coordinate (O) at ($(A)!{2/3}!(M)$);
			\coordinate (S) at ($(O)+(0,6)$);
			\tkzDrawPoints[fill=black](A,B,C,O,M,S)
			\tkzDrawSegments(S,A S,B S,C A,B B,C)
			\tkzDrawSegments[dashed](A,C S,O A,M)
			\tkzLabelPoints[above](S)
			\tkzLabelPoints[above left](A)
			\tkzLabelPoints[above right](C)
			\tkzLabelPoints[below](B,O)
			\tkzLabelPoints[below right](M)
			\end{tikzpicture}}}
\end{ex}

\begin{ex}%[GHK1 - Lương Thế Vinh CSA]%[Phan Anh - EX1]%[2H1B3-2]
	Cho lăng trụ đứng tam giác $ABC.A'B'C'$ có đáy $ABC$ là tam giác vuông cân tại $B$ với $BA=BC=a$, biết $A'B$ tạo với mặt phẳng $(ABC)$ một góc $60^\circ$. Thể tích $V$ của khối lăng trụ đã cho bằng
	\choice
	{\True $V=\dfrac{a^3\sqrt{3}}{2}$}
	{$V=2a^3$}
	{$V=\dfrac{a^3\sqrt{3}}{6}$}
	{$V=\dfrac{a^3}{2}$}
	\loigiai{\immini{Do $ABC.A'B'C'$ là lăng trụ đứng nên $AB$ là hình chiếu của $A'B$ lên $(ABC)$, suy ra 
			$$(A'B,(ABC))=(A'B,AB)=\widehat{A'BA}=60^\circ.$$
			Ta có $AA'=AB\tan60^\circ=a\sqrt{3}$. Nên thể tích lăng trụ là $$V=AA'\cdot S_{ABC}=a\sqrt{3}\cdot\dfrac{a^2}{2}=\dfrac{a^3\sqrt{3}}{2}.$$}
		{\begin{tikzpicture}[line cap=round, line join=round,font=\footnotesize,>=stealth, scale=0.6]
			\tkzDefPoints{0/0/A,1.5/-1.2/B,5/0/C,0/4/A'}
			\coordinate (B') at ($(A')+(B)-(A)$);
			\coordinate (C') at ($(A')+(C)-(A)$);
			\tkzDrawPoints[fill=black](A,B,C,A',B',C')
			\tkzDrawSegments(A,B B,C C,C' B,B' A,A' A',B' B',C' C',A' A',B)
			\tkzDrawSegments[dashed](A,C)
			\tkzLabelPoints[above left](A,A')
			\tkzLabelPoints[above right](C,C')
			\tkzLabelPoints[below right](B,B')
			\tkzMarkRightAngle(C,B,A)
			\end{tikzpicture}}}
\end{ex}

\begin{ex}%[GHK1 - Lương Thế Vinh CSA]%[Phan Anh - EX1]%[2H2K2-2]
	Cho tứ diện $ABCD$ đều có $AB=a\sqrt{2}$. Bán kính $R$ của mặt cầu chứa tất cả các đỉnh $A$, $B$, $C$, $D$ là số nào sau đây?
	\choice
	{\True $R=\dfrac{a\sqrt{3}}{2}$}
	{$R=\dfrac{3a}{2}$}
	{$R=\dfrac{3\sqrt{2}a}{2}$}
	{$R=a\sqrt{3}$}
	\loigiai{\immini{Gọi $O$ là tâm tam giác đều $BCD$, ta có $AO$ là trục đường tròn ngoại tiếp tam giác $BCD$.\\
			Gọi $M$ là trung điểm của $AB$, kẻ trung trực $Mt$ cua cạnh $AB$ trong tam giác $ABO$.\\
			Gọi $I$ là giao điểm của $Mt$ và $AO$, khi đó $I$ là tâm mặt cầu ngoại tiếp $ABCD$.\\
			Xét tam giác $ABO$, ta có $BO=\dfrac{2}{3}\cdot\dfrac{a\sqrt{2}\cdot\sqrt{3}}{2}=\dfrac{a\sqrt{6}}{3}$.\\
			Suy ra $AO=\sqrt{AB^2-BO^2}=\dfrac{2a\sqrt{3}}{3}$.\\
			Ta có $\triangle AMI\backsim\triangle AOB$, nên suy ra}
		{\begin{tikzpicture}[line cap=round, line join=round,font=\footnotesize,>=stealth, scale=0.7]
			\tkzDefPoints{0/0/B,2/-1.5/C,6/0/D}
			\coordinate (N) at ($(C)!0.5!(D)$);
			\coordinate (O) at ($(B)!{2/3}!(N)$);
			\coordinate (A) at ($(O)+(0,6)$);
			\coordinate (M) at ($(A)!0.5!(B)$);
			\coordinate (I) at ($(A)!0.7!(O)$);
			\coordinate (t) at ($(M)!1.3!(I)$);
			\tkzDrawPoints[fill=black](A,B,C,O,D,I,M)
			\tkzDrawSegments(B,A A,C A,D C,D B,C)
			\tkzDrawSegments[dashed](B,D A,O B,O M,t)
			\tkzLabelPoints[above](A)
			\tkzLabelPoints[above left](B,M)
			\tkzLabelPoints[above right](D,I)
			\tkzLabelPoints[below](C,O,t)
			\end{tikzpicture}}
		$$\dfrac{IA}{AB}=\dfrac{AM}{AO}\Leftrightarrow R=IA=\dfrac{AB\cdot AM}{AO}=\dfrac{a\sqrt{2}\cdot\dfrac{a\sqrt{2}}{2}}{\dfrac{2a\sqrt{3}}{3}}=\dfrac{a\sqrt{3}}{2}.$$}
\end{ex}
\begin{ex}%[GHK1 - Lương Thế Vinh CSA]%[Phan Anh - EX1]%[2D1K3-1]
	Có tất cả bao nhiêu giá trị của tham số $m$ để giá trị nhỏ nhất của hàm số $y=\dfrac{4x+m}{x+1}$ trên đoạn $[0;1]$ bằng $2$?
	\choice
	{\True $1$}
	{$0$}
	{$2$}
	{$3$}
	\loigiai{Ta có $y'=\dfrac{4-m}{(x+1)^2}$. Xét hai trường hợp
		\begin{itemize}
			\item \textbf{TH1:} $4-m>0\Leftrightarrow m<4$, khi đó ta có
			$$\min\limits_{x\in[0;1]}y=2\Leftrightarrow y(0)=2\Leftrightarrow m=2.\text{ (nhận)}$$
			\item \textbf{TH1:} $4-m<0\Leftrightarrow m>4$, khi đó ta có
			$$\min\limits_{x\in[0;1]}y=2\Leftrightarrow y(1)=2\Leftrightarrow\dfrac{m+4}{2}=2\Leftrightarrow m=0.\text{ (loại)}$$
		\end{itemize}
		Vậy chỉ có $1$ giá trị của $m$ để bài toán xảy ra.}
\end{ex}
\begin{ex}%[GHK1 - Lương Thế Vinh CSA]%[Phan Anh - EX1]%[2D1K5-1]
	\immini{Đường cong ở hình bên là đồ thị của hàm số $y=\dfrac{ax+b}{cx-2}$ với $a$, $b$, $c$ là các số thực khác $0$. Trong các số $a$, $b$, $c$ có bao nhiêu số dương?
		\choice
		{\True $3$}{$0$}{$1$}{$2$}}{\begin{tikzpicture}[scale=0.7, font=\footnotesize, line join=round, line cap=round, >=stealth,yscale=0.6]
		\draw[->](-4,0)--(6,0)node[above]{$x$};
		\draw[->](0,-4.7)--(0,6.8)node[right]{$y$};
		\draw[smooth,samples=100,domain=-4:1.3]plot(\x,{(\x+2)/(\x-2)});
		\draw[smooth,samples=100,domain=2.7:6]plot(\x,{(\x+2)/(\x-2)});
		\fill (0,0)node[below right]{$O$}circle(1.2pt) (0,1)node[above left]{$1$}circle(1.2pt) (2,0)node[below right]{$2$}circle(1.2pt);
		\draw(-4,1)--(6,1) (2,-4.7)--(2,6.8);
		\end{tikzpicture}}
	\loigiai{Dựa vào đồ thị ta có 
		\begin{itemize}
			\item $x=2$ là tiệm cận đứng nên $x=2$ là nghiệm của mẫu thức, suy ra $2x-2=0\Leftrightarrow c=1$.
			\item $y=1$ là tiệm cận ngang nên $\dfrac{a}{c}=1\Leftrightarrow a=1$.
			\item Ta có giao điểm của đồ thị với trục hoành là $\left(0;-\dfrac{b}{a}\right)$, ta có $-\dfrac{b}{a}<0\Leftrightarrow b>0$.
		\end{itemize}
		Vậy cả ba số $a$, $b$, $c$ đều dương.}
\end{ex}
\begin{ex}%[GHK1,Lương Thế Vinh -CS A - Hà Nội,2019 - 2020]%[Lê Quốc Hiệp, 12EX1]%[2H1B3-2]
	Cho hình lăng trụ $ABC.A'B'C'$ có đáy $ABC$ là tam giác đều cạnh $a$, độ dài cạnh bên là $a\sqrt{3}$. Biết rằng hình chiếu vuông góc của $A'$ lên $(ABC)$ là trung điểm $BC$. Thể tích của khối lăng trụ $ABC.A'B'C'$ là
	\choice
	{$\dfrac{2a^3\sqrt{3}}{3}$}
	{$\dfrac{a^3\sqrt{3}}{2}$}
	{$\dfrac{a^3\sqrt{3}}{8}$}
	{\True $\dfrac{3a^3\sqrt{3}}{8}$}
	\loigiai{
		\immini
		{Ta có $A'H$ là đường cao của lăng trụ $ABC.A'B'C'$.\\
			Xét tam giác vuông $A'HA$ có
			\[A'H=\sqrt{A'A^2-AH^2}=\sqrt{3a^2-\dfrac{3a^2}{4}}=\dfrac{3a}{2}.\]
			Vậy $V_{ABC.A'B'C'}=S_{ABC}\cdot A'H=\dfrac{a^2\sqrt{3}}{4}\cdot\dfrac{3a}{2}=\dfrac{3a^3\sqrt{3}}{8}$.
		}
		{\begin{tikzpicture}[line cap=round,line join=round,font=\footnotesize,>=stealth,scale=1]
			\fill (0,0) coordinate [label=above:$B$] (B) circle(1pt) 
			(2,0) coordinate [label=right:$A$] (A) circle(1pt)
			(-125:2) coordinate [label=below:$C$] (C) circle(1pt)
			($(B)!0.5!(C)$) coordinate [label=left:$H$] (H) circle(1pt)
			($(A)!0.5!(C)$) coordinate (M)
			($(H)!3!90:(M)$) coordinate [label=above:$A'$] (A') circle(1pt)
			($(A')-(A)$) coordinate [label=left:$B'$] (B') circle(1pt)
			($(C)+(B')$) coordinate [label=left:$C'$] (C') circle(1pt);
			\draw (A')--(C')--(B')--(A')--(A) (C')--(C)--(A);
			\draw[dashed] (C)--(B)--(A) (A')--(H)--(A) (B')--(B);
			\end{tikzpicture}}
	}
\end{ex}

\begin{ex}%[GHK1,Lương Thế Vinh -CS A - Hà Nội,2019 - 2020]%[Lê Quốc Hiệp, 12EX1]%[2H1B3-2]
	Cho hình chóp tam giác đều $S.ABC$ có độ dài cạnh đáy bằng $a$, góc giữa mặt bên và mặt đáy bằng $60^\circ$. Tính thể tích khối chóp đã cho
	\choice
	{$\dfrac{\sqrt{3}a^3}{6}$}
	{$\dfrac{\sqrt{3}a^3}{4}$}
	{\True $\dfrac{\sqrt{3}a^3}{24}$}
	{$\dfrac{\sqrt{3}a^3}{12}$}
	\loigiai{
		\immini
		{Gọi $M$ là trung điểm $BC$ và $G$ là trọng tâm tam giác $ABC$.\\
			Ta có $SG\perp(ABC)$ và $(SAG)$ vuông góc với $2$ mặt phẳng $(ABC)$, $(SBC)$. Suy ra $\widehat{SMG}=60^\circ$.\\
			Xét tam giác vuông $SGM$ có
			\[SG=GM\cdot\tan 60^\circ=\dfrac{a\sqrt{3}}{6}\cdot\sqrt{3}=\dfrac{a}{2}.\]
			Vậy $V_{S.ABC}=\dfrac{1}{3}SG\cdot S_{ABC}=\dfrac{1}{3}\cdot\dfrac{a^2\sqrt{3}}{4}\cdot\dfrac{a}{2}=\dfrac{\sqrt{3}a^3}{24}$.
		}
		{\begin{tikzpicture}[line cap=round,line join=round,font=\footnotesize,>=stealth,scale=1]
			\fill (0,0) coordinate [label=left:$A$] (A) circle(1pt) 
			(4,0) coordinate [label=right:$B$] (B) circle(1pt)
			(-30:4) coordinate [label=below:$C$] (C) circle(1pt)
			($(C)!0.5!(B)$) coordinate [label=right:$M$] (M) circle(1pt)
			($2/3*(M)$) coordinate [label=below:$G$] (G) circle(1pt)
			($(B)!2/3!(M)$) coordinate (x)
			($(G)!3!90:(x)$) coordinate [label=above:$S$] (S) circle(1pt);
			\draw (M)--(S)--(A)--(C)--(B)--(S)--(C);
			\draw [dashed] (M)--(A)--(B) (S)--(G);
			\tkzMarkAngle[size=.4](S,M,G)
			\end{tikzpicture}}
	}
\end{ex}

\begin{ex}%[GHK1,Lương Thế Vinh -CS A - Hà Nội,2019 - 2020]%[Lê Quốc Hiệp, 12EX1]%[2H1B3-2]
	Cho hình chóp $SABC$ có đáy là tam giác vuông tại $A$, $AC=a,BC=2a$, tam giác $SBC$ cân tại $S$ và nằm trong mặt phẳng vuông góc với đáy. Mặt bên $(SAC)$ hợp với mặt đáy một góc $60^\circ$. Thể tích của khối chóp $SABC$
	\choice
	{$\dfrac{a^3\sqrt{15}}{9}$}
	{$\dfrac{a^3\sqrt{3}}{5}$}
	{\True$\dfrac{a^3\sqrt{3}}{4}$}
	{$\dfrac{a^3\sqrt{3}}{12}$}
	\loigiai{
		\immini
		{Gọi $H$ là trung điểm $AB$, suy ra $SH\perp BC$ hay $SH\perp(ABC)$.\\
			Gọi $M$ là trung điểm $AC\Rightarrow HM\perp AC$ và $AC\perp SM$.\\
			Mà $AC$ là giao tuyến của $(SAC)$ và mặt đáy nên theo đề bài suy ra $\widehat{SMH}=60^\circ$.\\
			Ta có $HM=\dfrac{AB}{2}=\dfrac{\sqrt{BC^2-AC^2}}{2}=\dfrac{a\sqrt{3}}{2}$.\\
			Xét tam giác vuông $SHM$ có
			\[SH=MH\cdot\tan 60^\circ=\dfrac{a\sqrt{3}}{2}\cdot\sqrt{3}=\dfrac{3a}{2}.\]
		}
		{\begin{tikzpicture}[line cap=round,line join=round,font=\footnotesize,>=stealth,scale=1]
			\fill (0,0) coordinate [label=left:$B$] (B) circle(1pt) 
			(4,0) coordinate [label=right:$A$] (A) circle(1pt)
			(-125:2) coordinate [label=below:$C$] (C) circle(1pt)
			($(B)!0.5!(C)$) coordinate [label=below:$H$] (H) circle(1pt)
			($(A)!0.5!(C)$) coordinate [label=below:$M$] (M) circle(1pt)
			($(H)!2!90:(M)$) coordinate [label=left:$S$] (S) circle(1pt);
			\draw (S)--(C)--(A)--(S)--(M);
			\draw[dashed] (C)--(B)--(A) (M)--(H)--(S)--(B);
			\tkzMarkAngle[size=.5](S,M,H)
			\end{tikzpicture}}
		\noindent Suy ra $V_{SABC}=\dfrac{1}{3}SH\cdot S_{ABC}=\dfrac{1}{3}\cdot\dfrac{3a}{2}\cdot\dfrac{1}{2}\cdot a\sqrt{3}\cdot a=\dfrac{a^3\sqrt{3}}{4}$.
	}
\end{ex}

\begin{ex}%[GHK1,Lương Thế Vinh -CS A - Hà Nội,2019 - 2020]%[Lê Quốc Hiệp, 12EX1]%[2D1B1-3]
	Có tất cả bao nhiêu giá trị nguyên của tham số $m$ để hàm số $y=x^3+mx^2+2mx+2$ đồng biến trên $(-\infty;+\infty)$?
	\choice
	{\True $7$}
	{$5$}
	{vô số}
	{$6$}
	\loigiai{
		Hàm số đã cho có tập xác định $\mathscr{D}=\mathbb{R}$ và $y'=3x^2+2mx+2m$.\\
		Hàm số đã cho đồng biến trên $\mathbb{R}$ khi và chỉ khi
		\[y'\ge0,~\forall x\in\mathbb{R}\Leftrightarrow m^2-6m\le0\Leftrightarrow 0\le m\le6.\]
		Vậy $m\in\{0;1;2;3;4;5;6\}$.
	}
\end{ex}

\begin{ex}%[GHK1,Lương Thế Vinh -CS A - Hà Nội,2019 - 2020]%[Lê Quốc Hiệp, 12EX1]%[2D1K5-3]
	Có tất cả bao nhiêu giá trị nguyên của tham số $m$ để hàm số $y=|x^3-3x+m|$ có $5$ điểm cực trị? 
	\choice
	{$1$}
	{$2$}
	{\True $3$}
	{$4$}
	\loigiai{
		Xét hàm số $f(x)=x^3-3x+m\Rightarrow f'(x)=3x^2-3$, $f'(x)=0\Leftrightarrow x=\pm1$.\\
		Bảng biến thiên của hàm số $f(x)$ như sau
		\begin{center}
			\begin{tikzpicture}
			\tkzTabInit[nocadre=false,lgt=1.2,espcl=2.5,deltacl=0.6]
			{$x$ /.6,$f'(x)$ /.6,$f(x)$ /2}{$-\infty$,$-1$,$1$,$+\infty$}
			\tkzTabLine{,+,$0$,-,$0$,+,}
			\tkzTabVar{-/ $-\infty$ ,+/$2+m$,-/$-2+m$,+/$+\infty$}
			\end{tikzpicture}
		\end{center}
		Hàm số đã cho có $5$ điểm cực trị khi và chỉ khi đồ thị $y=f(x)$ cắt trục hoành tại $3$ điểm phân biệt.\\
		Dựa vào bảng biến thiên trên suy ra
		\[-2+m<0<2+m\Leftrightarrow -2<m<2.\]
		Vậy $m\in\{-1;0;1\}$.
	}
\end{ex}

\begin{ex}%[GHK1,Lương Thế Vinh -CS A - Hà Nội,2019 - 2020]%[Lê Quốc Hiệp, 12EX1]%[2D1K5-4]
	Biết rằng đường thẳng $y=x+m$ luôn cắt đồ thị hàm số $y=\dfrac{x+3}{x+1}$ tại hai điểm phân biệt, ký hiệu là $A,B$. Độ dài nhỏ nhất của đoạn thẳng $AB$ là
	\choice
	{$2\sqrt{2}$}
	{$2$}
	{\True $4$}
	{$\sqrt{2}$}
	\loigiai{
		Xét phương trình
		\[\dfrac{x+3}{x+1}=x+m\Leftrightarrow\heva{&x^2+mx+m-3=0~~(1)\\&x\ne-1.}\]
		Theo đề bài thì phương trình (1) luôn có $2$ nghiệm khác $-1$.\\
		Giả sử $A(x_1;x_1+m)$, $B(x_2;x_2+m)$, suy ra $AB^2=2(x_2-x_1)^2=2[(x_1+x_2)^2-4x_1x_2]$.\\
		Theo Vi-ét ta có $x_1+x_2=-m$ và $x_1x_2=m-3$, khi đó
		\[AB^2=2[m^2-4(m-3)]=2[(m-2)^2+8]\ge16.\tag{2}\]
		Từ (2) suy ra độ dài nhỏ nhất của đoạn thẳng $AB$ là $4$ khi $m=2$.
	}
\end{ex}

\begin{ex}%[GHK1,Lương Thế Vinh -CS A - Hà Nội,2019 - 2020]%[Lê Quốc Hiệp, 12EX1]%[2D1K5-4]
	Cho hàm số $f(x)=x^3-3x+2$. Có tất cả bao nhiêu giá trị nguyên của tham số $m$ thuộc đoạn $[-2019; 2019]$ để phương trình $f(|x|+m)=0$ có 4 nghiệm phân biệt? 
	\choice
	{\True $2017$}
	{$2019$}
	{$2018$}
	{$2020$}
	\loigiai{
		Ta có  $f(x)=0\Leftrightarrow x^3-3x+2=0\Leftrightarrow\hoac{&x=-2\\&x=1}$.\\
		Khi đó,\[f(|x|+m)=0\Leftrightarrow \hoac{&|x|+m=1\\&|x|+m=-2}\Leftrightarrow\hoac{&|x|=1-m\\&|x|=-2-m}.\tag{1}\]
		Từ (1) suy ra phương trình $f(|x|+m)=0$ có 4 nghiệm phân biệt khi và chỉ khi
		\[\heva{&1-m>0\\&-2-m>0}\Leftrightarrow \heva{&m<1\\&m<-2}\Leftrightarrow m<-2.\]
		Do $m$ nguyên và thuộc đoạn $[-2019; 2019]$ nên có $2017$ giá trị thỏa mãn đề bài.
	}
\end{ex}

\begin{ex}%[GHK1,Lương Thế Vinh -CS A - Hà Nội,2019 - 2020]%[Lê Quốc Hiệp, 12EX1]%[1H3K5-4]
	Cho hình chóp $SABCD$ có $ABCD$ là hình vuông cạnh $a$. Tam giác $SAB$ đều và nằm trong mặt phẳng vuông góc với mặt phẳng $(ABCD)$. Khoảng cách giữa hai đường thẳng $BC$ và $SD$ là
	\choice
	{$a$}
	{$\dfrac{a\sqrt{2}}{2}$}
	{$\dfrac{a\sqrt{3}}{3}$}
	{\True  $\dfrac{a\sqrt{3}}{2}$}
	\loigiai{
		\immini
		{Gọi $H$ là trung điểm $AB$ suy ra $SH\perp(ABCD)$.\\
			Gọi $K$ là trung điểm $SA\Rightarrow BK\perp SA$.\\
			Ta có $AD\perp AB$ và $AD\perp SH$ nên $AD\perp(SAB)$ suy ra $AD\perp BK$.\\
			Mà $BK\perp SA$ suy ra $BK\perp(SAB)$ hay
			\[\mathrm{d}(B,(SAD))=BK=\dfrac{a\sqrt{3}}{2}.\]
			Do $BC\parallel AD\Rightarrow BC\parallel (SAD)$ nên
			\[\mathrm{d}(BC,SD)=\mathrm{d}(BC,(SAD))=\mathrm{d}(B,(SAD))=\dfrac{a\sqrt{3}}{2}.\]
		}
		{\begin{tikzpicture}[line cap=round,line join=round,font=\footnotesize,>=stealth,scale=1]
			\fill (0,0) coordinate [label=below:$B$] (B) circle(1pt) 
			(4,0) coordinate [label=right:$C$] (C) circle(1pt)
			(-125:2) coordinate [label=below:$A$] (A) circle(1pt)
			($(B)!0.5!(A)$) coordinate [label=below:$H$] (H) circle(1pt)
			($(A)!0.5!(C)$) coordinate [label=below:$O$] (M) circle(1pt)
			($(H)!2!90:(M)$) coordinate [label=left:$S$] (S) circle(1pt)
			($(S)!0.5!(A)$) coordinate [label=left:$K$] (K) circle(1pt)
			($2*(M)$) coordinate [label=below:$D$] (D) circle(1pt)
			($(S)!2/3!(H)$) coordinate (G) node[shift={(-120:2ex)}]{$G$} circle(1pt);
			\draw (S)--(A)--(D)--(C)--(S)--(D);
			\draw[dashed] (A)--(B)--(C)--(A) (H)--(S)--(B)--(D)--(H) (B)--(K);
			\tkzMarkRightAngles[size=0.2](B,K,A)
			\end{tikzpicture}}
		
		
	}
\end{ex}

\begin{ex}%[GHK1,Lương Thế Vinh -CS A - Hà Nội,2019 - 2020]%[Lê Quốc Hiệp, 12EX1]%[2H2K2-2]
	Cho hình chóp $S.ABCD$ có đáy $ABCD$ là hình thoi cạnh $a,\widehat{ABC}=120^\circ,\triangle SAB$ đều và nằm trong mặt phẳng vuông góc với đáy. Tính bán kính mặt cầu ngoại tiếp hình chóp $S.ABC$.
	\choice
	{$\dfrac{a\sqrt{41}}{6}$}
	{\True $\dfrac{a\sqrt{39}}{6}$}
	{$\dfrac{a\sqrt{37}}{6}$}
	{$\dfrac{a\sqrt{35}}{6}$}
	\loigiai{
		\immini
		{Vì $ABCD$ là hình thoi có $\widehat{ABC}=120^\circ$ nên các tam giác $ABD$ và $DBC$ đều.\\
			Suy ra $D$ là tâm đường tròn ngoại tiếp tam giác $ABC$.\\\
			Gọi $H$ là trung điểm $AB$ suy ra $SH\perp(ABCD)$.\\
			Gọi $G$ là trọng tâm tam giác đều $SAB\Rightarrow G$ là tâm đường tròn ngoại tiếp tam giác $SAB$.\\
			Dựng các trục đường tròn $Dx$ của tam giác $ABC$ ($Dx\parallel SH$) và $Gy$ của tam giác $SAB$ ($Gy\parallel DH$). 
		}
		{\begin{tikzpicture}[line cap=round,line join=round,font=\footnotesize,>=stealth,scale=1]
			\fill (0,0) coordinate [label=below:$B$] (B) circle(1pt) 
			(4,0) coordinate [label=right:$C$] (C) circle(1pt)
			(-125:2) coordinate [label=below:$A$] (A) circle(1pt)
			($(B)!0.5!(A)$) coordinate [label=below:$H$] (H) circle(1pt)
			($(A)!0.5!(C)$) coordinate [label=below:$O$] (M) circle(1pt)
			($(H)!2!90:(M)$) coordinate [label=left:$S$] (S) circle(1pt)
			($2*(M)$) coordinate [label=below:$D$] (D) circle(1pt)
			($(S)!2/3!(H)$) coordinate (G) node[shift={(-120:2ex)}]{$G$} circle(1pt)
			($(D)+(S)-(H)$) coordinate [label=right:$x$] (x)
			($(D)+(G)-(H)$) coordinate [label=below right:$I$] (y) circle(1pt)
			($(G)!0.8!(y)$) coordinate (z)
			($(G)!1.4!(y)$) coordinate [label=below:$y$] (t);
			\draw (y)--(S)--(A)--(D)--(C)--(S)--(D)--(x) (z)--(t);
			\draw[dashed] (A)--(B)--(C)--(A) (H)--(S)--(B)--(D)--(H) (D) (G)--(z);
			\end{tikzpicture}}
		\noindent Gọi $I=Gy\cap Dx\Rightarrow IA=IB=IS=IC$, do đó $I$ là tâm mặt cầu $(S)$ ngoại tiếp hình chóp $S.ABC$.
		Bán kính của $(S)$ là
		\[R=IS=\sqrt{IG^2+SG^2}=\sqrt{DH^2+SG^2}=\sqrt{\dfrac{3a^2}{4}+\dfrac{a^2}{3}}=\dfrac{a\sqrt{39}}{6}.\]
	}
\end{ex}

\begin{ex}%[GHK1,Lương Thế Vinh -CS A - Hà Nội,2019 - 2020]%[Lê Quốc Hiệp, 12EX1]%[2D1K5-6]
	Có tất cả bao nhiêu tiếp tuyến của đồ thị hàm số $y=\dfrac{2x+1}{x-1}$ đi qua điểm $A(3;2)$?
	\choice
	{$0$}
	{$3$}
	{$2$}
	{\True $1$}
	\loigiai{
		Đường thẳng qua $A(3;2)$ có dạng $d\colon y=a(x-3)+2=ax-3a+2$.\\
		Xét phương trình
		\[\dfrac{2x+1}{x-1}=ax-3a+2\Leftrightarrow\heva{&ax^2-4ax+3a-3=0~~(1)\\&x\ne1.}\]
		Phương trình (1) luôn có nghiệm khác $-1$, suy ra $d$ tiếp xúc với đồ thị $y=\dfrac{2x+1}{x-1}$ khi và chỉ khi phương trình (1) có $1$ nghiệm.
		\begin{itemize}
			\item Với $a=0$, phương trình (1) vô nghiệm.
			\item Với $a\ne0$, phương trình (1) có $1$ nghiệm khi và chỉ khi
			\[\Leftrightarrow 4a^2-a(3a-3)=0\Leftrightarrow a^2+3a=0\Leftrightarrow \hoac{&a=0~\text{(loại)}\\&a=-3~\text{(nhận).}}\]
		\end{itemize}
		Vậy có duy nhất một tiếp tuyến đi qua $A(3;2)$.
	}
\end{ex}

\begin{ex}%[GHK1,Lương Thế Vinh -CS A - Hà Nội,2019 - 2020]%[Lê Quốc Hiệp, 12EX1]%[2H1K3-3]
	Cho hình hộp $ABCD.A'B'C'D'$ có thể tích $V$. Gọi $M,N$ là các điểm trên các cạnh $BB',DD'$ thỏa mãn $BM=\dfrac{1}{4}BB',DN=\dfrac{3}{4}DD'$. Tính thể tích khối tứ diện $ACMN$ theo $V$.
	\choice
	{$\dfrac{1}{3}V$}
	{$\dfrac{1}{12}V$}
	{$\dfrac{1}{4}V$}
	{\True $\dfrac{1}{6}V$}
	\loigiai{
		\immini
		{Gọi $O$ là tâm của đáy $ABCD$; $B,~h$ lần lượt là diện tích đáy và chiều cao của hình hộp $ABCD.A'B'C'D'$.\\
			Trong $(BDD'B')$, gọi $E=MN\cap BD\Rightarrow B$ là trung điểm $OE$.\\
			Ta có $S_{ACE}=2S_{ABC}=B$. Mặt khác,
			\begin{itemize}
				\item $V_{N.AECD}=\dfrac{1}{3}\mathrm{d}(N,(AECD))\cdot S_{AECD}=\dfrac{1}{3}\cdot\dfrac{3}{4}h\cdot\dfrac{3}{2}B=\dfrac{3}{8}V$.
				\item $V_{N.ACD}=\dfrac{1}{3}\mathrm{d}(N,(ACD))\cdot S_{ACD}=\dfrac{1}{3}\cdot\dfrac{3}{4}h\cdot\dfrac{1}{2}B=\dfrac{1}{8}V$.
				\item $V_{M.ACE}=\dfrac{1}{3}\mathrm{d}(M,(ACE))\cdot S_{ACE}=\dfrac{1}{3}\cdot\dfrac{1}{4}h\cdot B=\dfrac{1}{12}V$.
			\end{itemize}
		}
		{\begin{tikzpicture}[line cap=round,line join=round,font=\footnotesize,>=stealth,scale=0.8]
			\fill (0,0) coordinate [label=left:$A$] (A) circle(1pt)
			(4,0) coordinate (B) node[shift={(-90:1ex)}]{$B$} circle(1pt)
			(40:2) coordinate [label=left:$D$] (D) circle(1pt)
			(80:3) coordinate [label=left:$A'$] (A') circle(1pt)
			($(B)+(D)$) coordinate [label=right:$C$] (C) circle(1pt)
			($(A')+(B)$) coordinate [label=right:$B'$] (B') circle(1pt)
			($(A')+(C)$) coordinate [label=right:$C'$] (C') circle(1pt)
			($(A')+(D)$) coordinate [label=above:$D'$] (D') circle(1pt)
			($(B)!0.5!(D)$) coordinate [label=above:$O$] (O) circle(1pt)
			($(B)!1/4!(B')$) coordinate [label=below left:$M$] (M) circle(1pt)
			($(D)!3/2!(B)$) coordinate [label=below left:$E$] (E) circle(1pt)
			($(D)!3/4!(D')$) coordinate [label=left:$N$] (N) circle(1pt);
			\draw (C')--(B')--(A')--(A) (C)--(C')--(D')--(A') (B')--(B) (E)--(A)--(M)--(C)--(E)--(M);
			\draw [dashed] (A)--(D)--(D') (E)--(D)--(C) (N)--(M) (N)--(A)--(C)--(N) (A)--(B)--(C);
			\end{tikzpicture}}
		\noindent Khi đó, $V_{ACMN}=V_{N.AECD}-V_{N.ACD}-V_{M.ACE}=\dfrac{3}{8}V-\dfrac{1}{8}V-\dfrac{1}{12}V=\dfrac{1}{6}V$.
	}
\end{ex}

\begin{ex}%[GHK1,Lương Thế Vinh -CS A - Hà Nội,2019 - 2020]%[Lê Quốc Hiệp, 12EX1]%[2D1K5-3]
	Có tất cả bao nhiêu giá trị nguyên của tham số $m$ để phương trình $x^2+(2-m)x+9=0$ có nghiệm thuộc khoảng $(1;6)$? 
	\choice
	{$6$}
	{\True $4$}
	{$7$}
	{$5$}
	\loigiai{
		Vì $x=0$ không là nghiệm của phương trình đã cho nên chia cả 2 vế phương trình đã cho cho $x$ ta được phương trình
		\[x+\dfrac{9}{x}+2=m.\tag{1}\]
		Số nghiệm của phương trình (1) bằng số giao điểm của đồ thị hàm số $f(x)=x+\dfrac{9}{x}+2$ và $g(x)=m$.\\
		Ta có $f'(x)=1-\dfrac{9}{x^2}$, $f'(x)=0\Leftrightarrow 1-\dfrac{9}{x^2}=0\Leftrightarrow x=\pm3$.\\
		Lập bảng biến thiên của $f(x)$ trên $(1;6)$ ta được
		\begin{center}
			\begin{tikzpicture}
			\tkzTabInit[nocadre=false,lgt=1.2,espcl=2.5,deltacl=0.6]
			{$x$ /.6,$f(x)'$ /.6,$f(x)$ /2}{$1$,$3$,$6$}
			\tkzTabLine{,-,$0$,+,}
			\tkzTabVar{+/ $12$ ,-/$8$,+/$\dfrac{19}{2}$}
			\end{tikzpicture}
		\end{center}
		Vậy phương trình đã cho có nghiệm trên $(1;6)$ khi và chỉ khi $8\le m<12$, suy ra $m\in\{8;9;10;11\}$.
	}
\end{ex}

\begin{ex}%[GHK1,Lương Thế Vinh -CS A - Hà Nội,2019 - 2020]%[Lê Quốc Hiệp, 12EX1]%[2H1G3-3]
	Cho hình chóp $SABCD$ có đáy $ABCD$ là hình bình hành. Gọi $A',B',C',D'$ lần lượt là các điểm trên các cạnh $SA,SB,SC,SD$ sao cho $\dfrac{SA'}{SA}=\dfrac{SC'}{SC}=\dfrac{1}{2}, \dfrac{SB'}{SB}=\dfrac{SD'}{SD}=\dfrac{1}{3}$. Biết thể tích khối chóp $S.ABCD$ là $V$, tính thể tích tứ diện $A'B'C'D'$ theo $V$.
	\choice
	{$V_{A'B'C'D'}=\dfrac{V}{6}$}
	{$V_{A'B'C'D'}=\dfrac{V}{3}$}
	{\True $V_{A'B'C'D'}=\dfrac{V}{36}$}
	{$V_{A'B'C'D'}=\dfrac{V}{18}$}
	\loigiai{
		\immini
		{Theo đề bài ta có
			\[V=V_{A'B'C'.ABC}+V_{A'D'C'.ADC}+V_{S.B'D'C'}+V_{S.B'D'A'}+V_{A'B'C'D'}(*)\]
			Ta có
			$\dfrac{V_{S.A'B'C'}}{V_{S.ABC}}=\dfrac{SA'}{SA}\cdot\dfrac{SB'}{SB}\cdot\dfrac{SC'}{SC}=\dfrac{1}{12}$\\
			$\Rightarrow V_{S.A'B'C'}=\dfrac{1}{24}V$\\$\Rightarrow V_{A'B'C'.ABC}=\dfrac{1}{2}V-\dfrac{1}{24}V=\dfrac{11}{24}V$.
		}
		{\begin{tikzpicture}[line cap=round,line join=round,font=\footnotesize,>=stealth,scale=1]
			\fill (0,0) coordinate [label=below:$B$] (B) circle(1pt) 
			(4,0) coordinate [label=right:$C$] (C) circle(1pt)
			(-145:2) coordinate [label=below:$A$] (A) circle(1pt)
			($(B)!0.5!(A)$) coordinate (H)
			($(A)!0.5!(C)$) coordinate (M)
			($(H)!2!90:(M)$) coordinate [label=left:$S$] (S) circle(1pt)
			($(S)!0.5!(A)$) coordinate [label=left:$A'$] (A') circle(1pt)
			($2*(M)$) coordinate [label=below:$D$] (D) circle(1pt)
			($(S)!0.5!(C)$) coordinate [label=right:$C'$] (C')  circle(1pt)
			($(S)!1/3!(B)$) coordinate (B') node[shift={(180:1.2ex)}]{$B'$} circle(1pt)
			($(S)!1/3!(D)$) coordinate (D') node[shift={(-105:1.7ex)}]{$D'$} circle(1pt);
			\draw (S)--(A)--(D)--(C)--(S)--(D) (A')--(D')--(C');
			\draw[dashed] (A)--(B)--(C)--(A) (S)--(B)--(D) (A')--(B')--(C') (B')--(D') (A')--(C');
			\end{tikzpicture}}
		\noindent Tính tương tự như trên, suy ra
		\begin{itemize}
			\item $\dfrac{V_{S.A'D'C'}}{V_{S.ADC}}=\dfrac{SA'}{SA}\cdot\dfrac{SD'}{SD}\cdot\dfrac{SC'}{SC}=\dfrac{1}{12}
			\Rightarrow V_{A'D'C'.ADC}=\dfrac{11}{24}V$.
			\item $\dfrac{V_{S.B'D'C'}}{V_{S.BDC}}=\dfrac{SB'}{SB}\cdot\dfrac{SD'}{SD}\cdot\dfrac{SC'}{SC}=\dfrac{1}{18}
			\Rightarrow V_{S.B'D'C'}=\dfrac{1}{36}V$.
			\item $\dfrac{V_{S.B'D'A'}}{V_{S.BDA}}=\dfrac{SB'}{SB}\cdot\dfrac{SD'}{SD}\cdot\dfrac{SA'}{SA}=\dfrac{1}{18}
			\Rightarrow V_{S.B'D'A'}=\dfrac{1}{36}V$.
		\end{itemize}
		Từ $(*)$ suy ra 
		\[V_{A'B'C'D'}=V-\dfrac{11}{24}V-\dfrac{11}{24}V-\dfrac{1}{36}V-\dfrac{1}{36}V=\dfrac{1}{36}V.\]
	}
\end{ex}

\begin{ex}%[GHK1,Lương Thế Vinh -CS A - Hà Nội,2019 - 2020]%[Lê Quốc Hiệp, 12EX1]%[2D1G5-5]
	Cho $f(x)$ là đa thức bậc bốn của $x$. Đồ thị hàm số $f'(x)$ có điểm cực đại là $A(1;0)$, và đi qua các điểm $B(0;-4)$, $C(4;0)$. Hàm số $y=f(x^2+3x)$ nghịch biến trên khoảng nào dưới đây?	\choice
	{$(1;2)$}
	{$(-4;-3)$}
	{\True$(0;1)$}
	{$(-2;-1)$}
	\loigiai{
		Đặt $g(x)=f'(x)=ax^3+bx^2+cx+d$ ($a,b,c,d\in\mathbb{R}$).\\
		Theo đề bài ta có hệ phương trình
		\[\heva{&g'(1)=0\\&g(1)=0\\&g(0)=-4\\&g(4)=0}\Leftrightarrow\heva{&3a+2b+c=0\\&a+b+c+d=0\\&d=-4\\&64a+16b+4c+d=0}\Leftrightarrow\heva{&a=1\\&b=-6\\&c=9\\&d=-4.}\]
		Suy ra $f'(x)=x^3-6x^2+9x-4$, $f'(x)=0\Leftrightarrow x=1$ hoặc $x=4$.\\
		Bảng biến thiên của hàm số $f(x)$ như sau
		\begin{center}
			\begin{tikzpicture}
			\tkzTabInit[nocadre=false,lgt=1.2,espcl=2.5,deltacl=0.6]
			{$x$ /.6,$f'(x)$ /.6,$f(x)$ /2}{$-\infty$,$1$,$4$,$+\infty$}
			\tkzTabLine{,-,$0$,-,$0$,+,}
			\tkzTabVar{+/ $+\infty$ ,R,-/$f(4)$,+/$+\infty$}
			\end{tikzpicture}
		\end{center}
		Xét hàm số $y=f(x^2+3x)$, $y'=(2x+3)\cdot f'(x^2+3x)$.\\
		Ta có $f'(x^2+3x)>0\Leftrightarrow x^2+3x>4\Leftrightarrow x<-4\text{ hoặc }x>1$.\\
		Bảng xét dấu của $y'=(2x+3)f'(x^2+3x)$ như sau
		\begin{center}
			\begin{tikzpicture}
			\tkzTabInit[nocadre=false,lgt=2.5,espcl=2.5,deltacl=0.6]
			{$x$ /1,$2x+3$ /.9,$f'(x^2+3x)$ /.9,$y'$ /.9}
			{$-\infty$,$-4$,$-\dfrac{3}{2}$,$1$,$+\infty$}
			\tkzTabLine{,-,t,-,$0$,+,t,+}
			\tkzTabLine{,+,$0$,-,t,-,$0$,+}
			\tkzTabLine{,-,$0$,+,$0$,-,$0$,+}
			\end{tikzpicture}
		\end{center}
		Dựa vào bảng xét dấu trên và dựa vào các khoảng mà đề bài đưa ra, suy ra hàm số $y=f(x^2+3x)$ nghịch biến trên $(0;1)$.
	}
\end{ex}

\begin{ex}%[GHK1,Lương Thế Vinh -CS A - Hà Nội,2019 - 2020]%[Lê Quốc Hiệp, 12EX1]%[1H3G5-3]
	Cho hình chóp $SABC$ có đáy là tam giác $ABC$ vuông tại $A, \widehat{ABC}=30^\circ$, tam giác $SBC$ là tam giác đều cạnh $a$ và nằm trong mặt phẳng vuông góc với đáy. Khoảng cách từ $C$ đến $(SAB)$ là
	\choice
	{$\dfrac{a\sqrt{39}}{52}$}
	{\True $\dfrac{a\sqrt{39}}{13}$}
	{$\dfrac{a\sqrt{39}}{26}$}
	{$\dfrac{2a\sqrt{39}}{13}$}
	\loigiai{
		\immini
		{Gọi $H$ là trung điểm $BC\Rightarrow SH\perp BC\Rightarrow SH\perp(ABC)$.\\
			Gọi $K$ là trung điểm $AB$ và $E$ là hình chiếu vuông góc của $H$ lên $SK$.\\
			Ta có $AB\perp HK$ và $AB\perp SH$ nên $AB\perp (SHK)\Rightarrow AB\perp HE$.\\
			Khi đó, $HE\perp SK$ và $HE\perp AB$ nên $HE\perp(SAB)$, suy ra
			\[\mathrm{d}(C,(SAB))=2\mathrm{d}(H,(SAB))=2HE.\]
			Tam giác vuông $ABC$ có $AC=BC\cdot\sin30^\circ=\dfrac{a}{2}$\\$\Rightarrow HK=\dfrac{AC}{2}=\dfrac{a}{4}$.
		}
		{\begin{tikzpicture}[line cap=round,line join=round,font=\footnotesize,>=stealth,scale=1]
			\fill (0,0) coordinate [label=left:$B$] (B) circle(1pt) 
			(4,0) coordinate [label=right:$A$] (A) circle(1pt)
			(-25:3) coordinate [label=below:$C$] (C) circle(1pt)
			($(B)!0.5!(C)$) coordinate [label=below:$H$] (H) circle(1pt)
			($(B)!0.5!(A)$) coordinate [label=below:$K$] (K) circle(1pt)
			($(A)!0.5!(C)$) coordinate (M)
			($(H)!2!90:(M)$) coordinate [label=left:$S$] (S) circle(1pt)
			($(S)!0.7!(K)$) coordinate [label=left:$E$] (E) circle(1pt);
			\draw (C)--(B)--(S)--(C)--(A)--(S) (H)--(S);
			\draw[dashed] (B)--(A) (E)--(H)--(K) (S)--(K);
			\tkzMarkAngle[size=.6](C,B,A)
			\end{tikzpicture}}
		\noindent Tam giác vuông $SHK$ có
		\[\dfrac{1}{HE^2}=\dfrac{1}{SH^2}+\dfrac{1}{HK^2}=\dfrac{4}{3a^2}+\dfrac{16}{a^2}=\dfrac{52}{3a^2}\Rightarrow HE=\dfrac{a\sqrt{39}}{26}.\]
		Vậy $\mathrm{d}(C,(SAB))=2\cdot\dfrac{a\sqrt{39}}{26}=\dfrac{a\sqrt{39}}{13}$.
	}
\end{ex}
\Closesolutionfile{ans}
\begin{indapan}{10}
	{ans/ans-2-GHK1-8-GHK1-LuongTheVinh-CSA-HaNoi-20}
\end{indapan}